\section{Appendix:}
\textbf{Preliminary: LoRA}
Hu \textit{et al.} \cite{hu2022lora} propose that pre-trained model parameters can be efficiently updated using low-rank decomposition. In the context of FFN modules within Transformer blocks, consider the weight matrices $A$ and $B$. These matrices can be updated as:
\begin{equation}
A' = A + U_1V_1,
\end{equation}
where $U_1 \in \mathbb{R}^{d \times r}$ and $V_1 \in \mathbb{R}^{r \times d}$ are low-rank matrices with $r \ll d$. During training, only the LoRA matrices $U_1, V_1$ are trainable, while the original pre-trained weights $A$ remain frozen. This approach significantly reduces the number of trainable parameters while enabling effective knowledge modification in the given model.

\noindent
\textbf{Implementation Details}
Below are the hyperparameters used for pretraining and multi-step classification \ref{tab:impl-details-classification}.
\begin{table}[h!]
\centering
\begin{tabular}{l|c}
\textbf{Config} & \textbf{Value} \\
\hline
dataset           & CIFAR-100 \cite{krizhevsky2009learning} \\
backbone          & \textit{DeiT-Tiny Patch16} \cite{touvron2021training} \\
is pretrained     & True \\
learning rate     & 1e-4 \\
batch size        & 64 \\
training epochs   & 50 \\
optimizer         & AdamW \\
scheduler         & CosineAnnealingLR \\
number of classes & 100 \\
weight decay      & 0.01 \\
label smoothing   & 0.1 \\
alpha ($\alpha$)  & 1.0 \\
beta ($\beta$)  & 0.15 \\
gamma ($\gamma$)  & 0.15 \\
BND               & 4.7 \\
LoRA rank         & 8 \\
data percentage (retain) & 0.1 \\
data percentage (forget/new) & 1.0 \\
number of steps   & 5 \\
classes per step  & 10 \\
random seeds      & 42 \\
\end{tabular}
\caption{Training configuration used for Pretraining and BID-LoRA.}
\label{tab:impl-details-classification}
\end{table}


Below are the hyperparameters used for pretraining and multi-step face recognition.
\begin{table}[h!]
\centering
\begin{tabular}{l|c}
\textbf{Config} & \textbf{Value} \\
\hline
dataset           & CelebA-100 \cite{zhao2024continual}\\
backbone          & Face Transformer \cite{zhong2021face} \\
loss type         & ArcFace \\
is pretrained     & False \\
learning rate     & 1e-4 \\
head learning rate & 5e-4 \\
batch size        & 64 \\
training epochs   & 300 \\
optimizer         & AdamW \\
scheduler         & CosineAnnealingLR \\
number of classes & 100 \\
weight decay      & 1e-4 \\
gradient clipping & 1.0 \\
alpha ($\alpha$)  & 2.0 \\
beta ($\beta$)  & 0.05 \\
gamma ($\gamma$)  & 2.0 \\
BND               & 8.0 \\
LoRA rank         & 8 \\
ArcFace scale     & 64.0 \\
data percentage (retain) & 0.1 \\
data percentage (forget) & 1.0 \\
data percentage (new) & 1.0 \\
number of steps   & 5 \\
classes per step  & 10 \\
random seeds      & 42 \\
\end{tabular}
\caption{Training configuration used for Face Recognition Pretraining and BID-LoRA.}
\label{tab:impl-details-face}
\end{table}


\noindent
\textbf{Algorithmic View of BID-LoRA Framework}

Algorithm~\ref{algo:apply_lora} is used to apply BID-LoRA to a given model, Algorithm~\ref{algo:single_step} is carefully designed for implementing single-step adaptations, while Algorithm~\ref{algo:multi_step} applies multiple sequential steps for continual learning-unlearning.

\begin{algorithm}[htbp]
\caption{\textit{apply\_lora()}}
\label{algo:apply_lora}
\begin{algorithmic}[1]
\REQUIRE Pre-trained model $\mathcal{M}$ and rank $r$
\ENSURE Model with LoRA adapters $\mathcal{M}_{\text{LoRA}}$

\FOR{each layer in $\mathcal{M}$}
    \IF{layer name contains ``ffn'', ``mlp'', ``fc1'', or ``fc2''}
        \STATE $d_{\text{in}} \leftarrow \text{layer.in\_features}$
        \STATE $d_{\text{out}} \leftarrow \text{layer.out\_features}$
        \STATE $\text{LoRA}_A \leftarrow (r, d_{\text{in}}) \times 0.01$
        \STATE $\text{LoRA}_B \leftarrow (d_{\text{out}}, r) \times 0.00$
        \STATE Register $\text{LoRA}_A$ and $\text{LoRA}_B$ as trainable parameters
    \ENDIF
\ENDFOR
\STATE Freeze all parameters except classification head and LoRA parameters
\RETURN $\mathcal{M}_{\text{LoRA}}$
\end{algorithmic}
\end{algorithm}

\begin{algorithm}[htbp]
\caption{\textit{single\_step()}}
\label{algo:single_step}
\begin{algorithmic}[1]
\REQUIRE Model $\mathcal{M}$, datasets $D_f$, $D_r$, $D_{\text{new}}$, hyperparameters $\alpha$, $\beta$, $\gamma$, rank $r$
\ENSURE Updated model $\mathcal{M}_{\text{updated}}$

\STATE $\mathcal{M}_{\text{LoRA}} \leftarrow \textit{apply\_lora}(\mathcal{M}, r)$
\FOR{each training iteration}
    \STATE Sample batches: $(f_d, f_t) \sim D_f$, $(r_d, r_t) \sim D_r$, $(n_d, n_t) \sim D_{\text{new}}$
    \STATE $f_{\text{logits}} \leftarrow \mathcal{M}_{\text{LoRA}}(f_d)$
    \STATE $r_{\text{logits}} \leftarrow \mathcal{M}_{\text{LoRA}}(r_d)$
    \STATE $n_{\text{logits}} \leftarrow \mathcal{M}_{\text{LoRA}}(n_d)$
    
    \STATE $\mathcal{L}_{\text{forget}} \leftarrow \text{ReLU}(\text{BND} - \mathcal{L}(f_{\text{logits}}, f_t))$
    \STATE $\mathcal{L}_{\text{retain}} \leftarrow \mathcal{L}(r_{\text{logits}}, r_t)$
    \STATE $\mathcal{L}_{\text{new}} \leftarrow \mathcal{L}(n_{\text{logits}}, n_t)$
    
    \STATE $\mathcal{L}_{\text{total}} \leftarrow \alpha \mathcal{L}_{\text{retain}} + \beta \mathcal{L}_{\text{forget}} + \gamma \mathcal{L}_{\text{new}}$
    \STATE Update $\mathcal{M}_{\text{LoRA}}$ using $\nabla \mathcal{L}_{\text{total}}$
\ENDFOR
\RETURN $\mathcal{M}_{\text{updated}}$
\end{algorithmic}
\end{algorithm}


\begin{algorithm}[htbp]
\caption{\textit{multi\_step()}}
\label{algo:multi_step}
\begin{algorithmic}[1]
\REQUIRE Pre-trained model $\mathcal{M}_0$, complete dataset $D$, task sequence $T$
\ENSURE Final adapted model $\mathcal{M}_T$

\STATE $\mathcal{M}_{\text{current}} \leftarrow \mathcal{M}_0$
\FOR{each step $t \in \{1, 2, \ldots, T\}$}
    \STATE Generate task-specific datasets: $D_{f_t}$, $D_{r_t}$, $D_{\text{new}_t}$ from $D$
    \STATE Set hyperparameters: rank $r$, weights $\alpha$, $\beta$, $\gamma$
    \STATE $\mathcal{M}_{\text{current}} \leftarrow \textit{single\_step}(\mathcal{M}_{\text{current}}, D_{f_t}, D_{r_t}, D_{\text{new}_t}, \alpha, \beta, \gamma, r)$
\ENDFOR
\RETURN $\mathcal{M}_T \leftarrow \mathcal{M}_{\text{current}}$
\end{algorithmic}
\end{algorithm}